\title{Linguistik für Lehramtsstudierende}
\BackBody{\textit{Linguistik für Lehramtsstudierende} bietet eine grammatische Beschreibung der deutschen Standardsprache und zahlreicher Variationsphänomene. Die Darstellung setzt bei sprachlichen Formen an und untersucht deren Funktionen. Die Betrachtung bewegt sich von der größeren Ebene des Satzes über kleinere sprachliche Bausteine (Wortgruppen, Wörter, Wortbausteine, Silben) bis hin zu den kleinsten lautlichen Ebenen (Lautklassen und einzelne Laute). Den Abschluss bildet ein Kapitel zur Schreibung (Graphematik) des Deutschen. Auf allen sprachlichen Ebenen stellen wir zahlreiche Bezüge zu Schule und Vermittlung her. Begrifflich orientieren wir uns an dem aktuellen, von der Kultusministerkonferenz zustimmend zur Kenntnis genommenen \textit{Verzeichnis grundlegender grammatischer Fachausdrücke}.
In erster Linie ist dieses Lehrbuch als vertiefendes Begleitmaterial zu ein- und weiterführenden linguistischen Lehrveranstaltungen konzipiert. Es vermittelt vor allem Techniken der linguistischen Analyse und Reflexion. Viele Übungsaufgaben samt ausführlich kommentierter Lösungen ermöglichen auch den Einsatz im Selbststudium und zur berufsbegleitenden Weiterbildung. 
}
%\dedication{Change dedication in localmetadata.tex}
\typesetter{Marc Felfe, Jana Brunner}
\proofreader{Alena Witzlack,
Annika Schiefner,
Bruno Behling,
Frederic Blum,
Hannah Schleupner,
Hella Olbertz,
Jean Nitzke,
Katja Politt,
Leonie Twente,
Lisa Schäfer,
Ludger Paschen,
Niklas Reinken,
Rainer Schulze,
Sebastian Nordhoff,
Viktor Koehlich}
\author{Marc Felfe and Jana Brunner}
\BookDOI{10.5281/zenodo.21918884}
\renewcommand{\lsISBNdigital}{978-3-96110-593-9}
\renewcommand{\lsISBNsoftcover}{978-3-98554-210-9}
\renewcommand{\lsSeries}{tbls} % use lowercase acronym, e.g. sidl, eotms, tgdi
\renewcommand{\lsSeriesNumber}{22} %will be assigned when the book enters the proofreading stage
\renewcommand{\lsID}{596} % use lowercase acronym, e.g. sidl, eotms, tgdi


 \lsCoverTitleSizes{44pt}{15mm}



